% =====================================================================
%  fig_occlusion_3d.tex — 烟幕遮蔽几何示意图（对标 A196 图5-7）
% ---------------------------------------------------------------------
%  展示：导弹 M 看向真目标圆柱，视线锥；烟幕球 S 遮蔽视线；关键轮廓红高亮。
%  体现"两底面圆周精确归约"：圆柱全遮蔽 ⇔ 只查上/下底面圆周即可。
%  编译：cd templates/tikz && tectonic fig_occlusion_3d.tex
% =====================================================================
\documentclass[border=8pt]{standalone}
\usepackage{ctex}
% =====================================================================
%  tikz_geometry_preamble.tex — 几何/机理示意图 TikZ 公共导言（被 \input 的片段）
%  ---------------------------------------------------------------------
%  用法：主文档里 \documentclass[border=6pt]{standalone} + \usepackage{ctex,tikz}
%        之后 % =====================================================================
%  tikz_geometry_preamble.tex — 几何/机理示意图 TikZ 公共导言（被 \input 的片段）
%  ---------------------------------------------------------------------
%  用法：主文档里 \documentclass[border=6pt]{standalone} + \usepackage{ctex,tikz}
%        之后 % =====================================================================
%  tikz_geometry_preamble.tex — 几何/机理示意图 TikZ 公共导言（被 \input 的片段）
%  ---------------------------------------------------------------------
%  用法：主文档里 \documentclass[border=6pt]{standalone} + \usepackage{ctex,tikz}
%        之后 \input{tikz_geometry_preamble.tex}（本文件只含宏包/颜色/tikzset，无 documentclass）。
%  对标 A196 图5-8 的手法：虚线=隐藏边、灰渐变=实体、红粗线=关键轮廓、
%  标准几何符号标注点/线/圆。极简、黑白+关键色高亮，学术矢量风。
% =====================================================================
\usepackage{tikz}
\usepackage{amsmath,amssymb}
\usetikzlibrary{arrows.meta,calc,3d,shapes.geometric,shapes.misc,positioning,decorations.pathreplacing}

% ---- A196 风格配色（极简黑白 + 关键色红/蓝点缀）----
\definecolor{ink}{RGB}{40,40,40}        % 主墨色（非纯黑）
\definecolor{keyred}{RGB}{200,40,40}    % 关键轮廓/重点红
\definecolor{keyblue}{RGB}{60,110,180}  % 辅助蓝
\definecolor{fillsolid}{RGB}{235,235,235} % 实心浅灰
\definecolor{fillgrad}{RGB}{200,200,200}  % 渐变中灰(球体阴影)

% ---- 常用样式 定义（供图内复用）----
\tikzset{
  hidden/.style={draw=ink, dashed, line width=0.6pt},
  solid/.style={draw=ink, line width=0.6pt},
  keyline/.style={draw=keyred, line width=1.4pt},
  solidbody/.style={draw=ink, fill=fillsolid, line width=0.6pt},
  sphere/.style={draw=ink, fill=fillsolid, line width=0.6pt},
  proc/.style={draw=ink, rounded corners, fill=white, align=center, inner sep=5pt, font=\small},
  decision/.style={draw=ink, diamond, aspect=1.8, align=center, inner sep=3pt, font=\small},
  terminator/.style={draw=ink, rounded corners=10pt, fill=white, align=center, inner sep=5pt, font=\small},
  io/.style={draw=ink, trapezium, trapezium left angle=70, trapezium right angle=110, align=center, inner sep=4pt, font=\small},
  flow/.style={-{Stealth[length=2.6mm]}, thick, draw=ink},
  lbl/.style={font=\small, inner sep=1pt, fill=white},
}
（本文件只含宏包/颜色/tikzset，无 documentclass）。
%  对标 A196 图5-8 的手法：虚线=隐藏边、灰渐变=实体、红粗线=关键轮廓、
%  标准几何符号标注点/线/圆。极简、黑白+关键色高亮，学术矢量风。
% =====================================================================
\usepackage{tikz}
\usepackage{amsmath,amssymb}
\usetikzlibrary{arrows.meta,calc,3d,shapes.geometric,shapes.misc,positioning,decorations.pathreplacing}

% ---- A196 风格配色（极简黑白 + 关键色红/蓝点缀）----
\definecolor{ink}{RGB}{40,40,40}        % 主墨色（非纯黑）
\definecolor{keyred}{RGB}{200,40,40}    % 关键轮廓/重点红
\definecolor{keyblue}{RGB}{60,110,180}  % 辅助蓝
\definecolor{fillsolid}{RGB}{235,235,235} % 实心浅灰
\definecolor{fillgrad}{RGB}{200,200,200}  % 渐变中灰(球体阴影)

% ---- 常用样式 定义（供图内复用）----
\tikzset{
  hidden/.style={draw=ink, dashed, line width=0.6pt},
  solid/.style={draw=ink, line width=0.6pt},
  keyline/.style={draw=keyred, line width=1.4pt},
  solidbody/.style={draw=ink, fill=fillsolid, line width=0.6pt},
  sphere/.style={draw=ink, fill=fillsolid, line width=0.6pt},
  proc/.style={draw=ink, rounded corners, fill=white, align=center, inner sep=5pt, font=\small},
  decision/.style={draw=ink, diamond, aspect=1.8, align=center, inner sep=3pt, font=\small},
  terminator/.style={draw=ink, rounded corners=10pt, fill=white, align=center, inner sep=5pt, font=\small},
  io/.style={draw=ink, trapezium, trapezium left angle=70, trapezium right angle=110, align=center, inner sep=4pt, font=\small},
  flow/.style={-{Stealth[length=2.6mm]}, thick, draw=ink},
  lbl/.style={font=\small, inner sep=1pt, fill=white},
}
（本文件只含宏包/颜色/tikzset，无 documentclass）。
%  对标 A196 图5-8 的手法：虚线=隐藏边、灰渐变=实体、红粗线=关键轮廓、
%  标准几何符号标注点/线/圆。极简、黑白+关键色高亮，学术矢量风。
% =====================================================================
\usepackage{tikz}
\usepackage{amsmath,amssymb}
\usetikzlibrary{arrows.meta,calc,3d,shapes.geometric,shapes.misc,positioning,decorations.pathreplacing}

% ---- A196 风格配色（极简黑白 + 关键色红/蓝点缀）----
\definecolor{ink}{RGB}{40,40,40}        % 主墨色（非纯黑）
\definecolor{keyred}{RGB}{200,40,40}    % 关键轮廓/重点红
\definecolor{keyblue}{RGB}{60,110,180}  % 辅助蓝
\definecolor{fillsolid}{RGB}{235,235,235} % 实心浅灰
\definecolor{fillgrad}{RGB}{200,200,200}  % 渐变中灰(球体阴影)

% ---- 常用样式 定义（供图内复用）----
\tikzset{
  hidden/.style={draw=ink, dashed, line width=0.6pt},
  solid/.style={draw=ink, line width=0.6pt},
  keyline/.style={draw=keyred, line width=1.4pt},
  solidbody/.style={draw=ink, fill=fillsolid, line width=0.6pt},
  sphere/.style={draw=ink, fill=fillsolid, line width=0.6pt},
  proc/.style={draw=ink, rounded corners, fill=white, align=center, inner sep=5pt, font=\small},
  decision/.style={draw=ink, diamond, aspect=1.8, align=center, inner sep=3pt, font=\small},
  terminator/.style={draw=ink, rounded corners=10pt, fill=white, align=center, inner sep=5pt, font=\small},
  io/.style={draw=ink, trapezium, trapezium left angle=70, trapezium right angle=110, align=center, inner sep=4pt, font=\small},
  flow/.style={-{Stealth[length=2.6mm]}, thick, draw=ink},
  lbl/.style={font=\small, inner sep=1pt, fill=white},
}


\begin{document}

\begin{tikzpicture}[scale=1.0]
  % ---- 坐标系（3D 等轴测视角示意）----
  \coordinate (O) at (0,0);
  \draw[-{Stealth[length=2.2mm]}, ink] (O) -- (5,0) node[below, lbl]{$x$};
  \draw[-{Stealth[length=2.2mm]}, ink] (O) -- (0,2.6) node[left, lbl]{$z$};
  \draw[-{Stealth[length=2.2mm]}, ink] (O) -- (2.4,1.6) node[below right, lbl]{$y$};

  % ---- 导弹 M（进近方向来自 -x，位于左上方）----
  \coordinate (M) at (-1.0,2.3);
  \node[circle, fill=ink, inner sep=2.2pt] at (M) {};
  \node[lbl, above right] at (M) {$M$ 导弹};

  % ---- 真目标圆柱（3D 线段逼近, 灰实心侧影, 虚线=隐藏边）----
  \def\rr{1.0}
  \def\hh{1.8}
  % 圆柱侧面（下底→上底的两条竖直母线, 前实后虚）
  \draw[solid] (-\rr,0) -- ++(0,\hh);
  \draw[solid] (\rr,0) -- ++(0,\hh);
  \draw[hidden] (0,0) -- ++(0,\hh);          % 背面母线(虚线)
  % 下底面椭圆（前实后虚）
  \draw[solid] (-\rr,0) arc (180:360:{\rr} and {0.35});
  \draw[hidden] (0,0) arc (0:180:{\rr} and {0.35});
  \draw[solid] (\rr,0) arc (0:180:{\rr} and {0.35});
  % 上底面椭圆
  \draw[solid] (-\rr,\hh) arc (180:360:{\rr} and {0.35});
  \draw[hidden] (0,\hh) arc (0:180:{\rr} and {0.35});
  \draw[solid] (\rr,\hh) arc (0:180:{\rr} and {0.35});
  \node[below, lbl] at (0.5,-0.35) {真目标圆柱};
  \node[lbl] at (\rr+0.32, \hh/2) {$r_t$=7m};

  % ---- 关键轮廓红高亮：导弹视角的圆柱外包络（两底面圆周）----
  \draw[keyline] (M) -- (\rr, \hh);   % 视线到底缘
  \draw[keyline] (M) -- (\rr, 0);     % 视线到上缘
  \draw[keyline] (-\rr,\hh) arc (0:180:{\rr} and {0.35});  % 上底面圆周(关键)
  \draw[keyline] (\rr,\hh) arc (0:180:{\rr} and {0.35});
  \draw[keyline] (\rr,0) arc (0:180:{\rr} and {0.35});     % 下底面圆周(关键)

  % ---- 烟幕球 S（立体感: 灰球 + 高光, 位于导弹-目标连线上）----
  \node[lbl, align=center, font=\footnotesize] at (2.2,1.9)
    {$S$: 烟幕球 $R$=10m};
  \begin{scope}
    \clip (1.9,1.0) circle (0.7);
    \fill[fillgrad] (-3.2,1.0) circle (0.75);   % 用更大圆做球体灰度背景
    \fill[fillsolid] (1.9,1.0) circle (0.7);
    \fill[white] (1.5,1.45) circle (0.28);       % 高光
  \end{scope}
  \draw[solid] (1.9,1.0) circle (0.7);
  \draw[hidden] (1.9,1.0) ellipse ({0.7} and {0.22}); % 球上纬线(示意立体)
  \node[lbl, below] at (1.9,0.15) {$C$(球心)};

  % ---- 两底面圆周归约示意（右侧, 简洁）----
  \begin{scope}[shift={(4.6,-0.4)}, scale=0.85]
    \node[align=center, font=\footnotesize] at (0,1.9) {归约: 全遮蔽};
    \node[align=center, font=\footnotesize] at (0,1.5)
      {$\Leftrightarrow$ 上底$\cup$下底圆周};
    \draw[keyline] (0,0.3) circle (0.55);
    \draw[keyline] (0,-0.55) ellipse ({0.55} and {0.22});
    \node[lbl, font=\footnotesize] at (0.8,0.4) {顶圆周};
    \node[lbl, font=\footnotesize] at (0.8,-0.55) {底圆周};
    \node[align=center, font=\scriptsize, lbl] at (0.4,-1.4)
      {不需对整圆柱采样\\只查两圆周};
  \end{scope}

  % ---- 图注（放底部, 避开主体）----
  \node[font=\scriptsize, lbl] at (2.0,-2.0)
    {图：视线锥全遮蔽判定——圆柱由两底面圆周精确归约};
\end{tikzpicture}

\end{document}
